
\documentclass[a4paper,11pt]{article}
\usepackage[a4paper, margin=8em]{geometry}

% usa i pacchetti per la scrittura in italiano
\usepackage[french,italian]{babel}
\usepackage[T1]{fontenc}
\usepackage[utf8]{inputenc}
\frenchspacing 

% usa i pacchetti per la formattazione matematica
\usepackage{amsmath, amssymb, amsthm, amsfonts}

% usa altri pacchetti
\usepackage{gensymb}
\usepackage{hyperref}
\usepackage{standalone}

\usepackage{colortbl}

\usepackage{xstring}
\usepackage{karnaugh-map}

% imposta il titolo
\title{Embedded Operating Systems and Architectures notes}
\author{Luca Seggiani}
\date{2026}

% imposta lo stile
% usa helvetica
\usepackage[scaled]{helvet}
% usa palatino
\usepackage{palatino}
% usa un font monospazio guardabile
\usepackage{lmodern}

\renewcommand{\rmdefault}{ppl}
\renewcommand{\sfdefault}{phv}
\renewcommand{\ttdefault}{lmtt}

% circuiti
\usepackage{circuitikz}
\usetikzlibrary{babel}

% testo cerchiato
\newcommand*\circled[1]{\tikz[baseline=(char.base)]{
            \node[shape=circle,draw,inner sep=2pt] (char) {#1};}}

% disponi il titolo
\makeatletter
\renewcommand{\maketitle} {
	\begin{center} 
		\begin{minipage}[t]{.8\textwidth}
			\textsf{\huge\bfseries \@title} 
		\end{minipage}%
		\begin{minipage}[t]{.2\textwidth}
			\raggedleft \vspace{-1.65em}
			\textsf{\small \@author} \vfill
			\textsf{\small \@date}
		\end{minipage}
		\par
	\end{center}

	\thispagestyle{empty}
	\pagestyle{fancy}
}
\makeatother

% disponi teoremi
\usepackage{tcolorbox}
\newtcolorbox[auto counter, number within=section]{theorem}[2][]{%
	colback=blue!10, 
	colframe=blue!40!black, 
	sharp corners=northwest,
	fonttitle=\sffamily\bfseries, 
	title=Theorem~\thetcbcounter: #2, 
	#1
}

% disponi definizioni
\newtcolorbox[auto counter, number within=section]{definition}[2][]{%
	colback=red!10,
	colframe=red!40!black,
	sharp corners=northwest,
	fonttitle=\sffamily\bfseries,
	title=Definition~\thetcbcounter: #2,
	#1
}

% disponi codice
\usepackage{listings}
\usepackage[table]{xcolor}

\definecolor{codegreen}{rgb}{0,0.6,0}
\definecolor{codegray}{rgb}{0.5,0.5,0.5}
\definecolor{codepurple}{rgb}{0.58,0,0.82}
\definecolor{backcolour}{rgb}{0.95,0.95,0.92}

\lstdefinestyle{codestyle}{
		backgroundcolor=\color{black!5}, 
		commentstyle=\color{codegreen},
		keywordstyle=\bfseries\color{magenta},
		numberstyle=\sffamily\tiny\color{black!60},
		stringstyle=\color{green!50!black},
		basicstyle=\ttfamily\footnotesize,
		breakatwhitespace=false,         
		breaklines=true,                 
		captionpos=b,                    
		keepspaces=true,                 
		numbers=left,                    
		numbersep=5pt,                  
		showspaces=false,                
		showstringspaces=false,
		showtabs=false,                  
		tabsize=2
}

\lstdefinestyle{shellstyle}{
		backgroundcolor=\color{black!5}, 
		basicstyle=\ttfamily\footnotesize\color{black}, 
		commentstyle=\color{black}, 
		keywordstyle=\color{black},
		numberstyle=\color{black!5},
		stringstyle=\color{black}, 
		showspaces=false,
		showstringspaces=false, 
		showtabs=false, 
		tabsize=2, 
		numbers=none, 
		breaklines=true
}

\lstdefinelanguage{x86}{ 
  keywords={AAA, AAD, AAM, AAS, ADC, ADCB, ADCW, ADCL, ADD, ADDB, ADDW, ADDL, AND, ANDB, ANDW, ANDL,
        ARPL, BOUND, BSF, BSFL, BSFW, BSR, BSRL, BSRW, BSWAP, BT, BTC, BTCB, BTCW, BTCL, BTR, 
        BTRB, BTRW, BTRL, BTS, BTSB, BTSW, BTSL, CALL, CBW, CDQ, CLC, CLD, CLI, CLTS, CMC, CMP,
        CMPB, CMPW, CMPL, CMPS, CMPSB, CMPSD, CMPSW, CMPXCHG, CMPXCHGB, CMPXCHGW, CMPXCHGL,
        CMPXCHG8B, CPUID, CWDE, DAA, DAS, DEC, DECB, DECW, DECL, DIV, DIVB, DIVW, DIVL, ENTER,
        HLT, IDIV, IDIVB, IDIVW, IDIVL, IMUL, IMULB, IMULW, IMULL, IN, INB, INW, INL, INC, INCB,
        INCW, INCL, INS, INSB, INSD, INSW, INT, INT3, INTO, INVD, INVLPG, IRET, IRETD, JA, JAE,
        JB, JBE, JC, JCXZ, JE, JECXZ, JG, JGE, JL, JLE, JMP, JNA, JNAE, JNB, JNBE, JNC, JNE, JNG,
        JNGE, JNL, JNLE, JNO, JNP, JNS, JNZ, JO, JP, JPE, JPO, JS, JZ, LAHF, LAR, LCALL, LDS,
        LEA, LEAVE, LES, LFS, LGDT, LGS, LIDT, LMSW, LOCK, LODSB, LODSD, LODSW, LOOP, LOOPE,
        LOOPNE, LSL, LSS, LTR, MOV, MOVB, MOVW, MOVL, MOVSB, MOVSD, MOVSW, MOVSX, MOVSXB,
        MOVSXW, MOVSXL, MOVZX, MOVZXB, MOVZXW, MOVZXL, MUL, MULB, MULW, MULL, NEG, NEGB, NEGW,
        NEGL, NOP, NOT, NOTB, NOTW, NOTL, OR, ORB, ORW, ORL, OUT, OUTB, OUTW, OUTL, OUTSB, OUTSD,
        OUTSW, POP, POPL, POPW, POPB, POPA, POPAD, POPF, POPFD, PUSH, PUSHL, PUSHW, PUSHB, PUSHA, 
        RCLL, RCR, RCRB, RCRW, RCRL, RDMSR, RDPMC, RDTSC, REP, REPE, REPNE, RET, ROL, ROLB, ROLW,
        ROLL, ROR, RORB, RORW, RORL, SAHF, SAL, SALB, SALW, SALL, SAR, SARB, SARW, SARL, SBB,
        SBBB, SBBW, SBBL, SCASB, SCASD, SCASW, SETA, SETAE, SETB, SETBE, SETC, SETE, SETG, SETGE,
        SETL, SETLE, SETNA, SETNAE, SETNB, SETNBE, SETNC, SETNE, SETNG, SETNGE, SETNL, SETNLE,
        SETNO, SETNP, SETNS, SETNZ, SETO, SETP, SETPE, SETPO, SETS, SETZ, SGDT, SHL, SHLB, SHLW,
        SHLL, SHLD, SHR, SHRB, SHRW, SHRL, SHRD, SIDT, SLDT, SMSW, STC, STD, STI, STOSB, STOSD,
        STOSW, STR, SUB, SUBB, SUBW, SUBL, TEST, TESTB, TESTW, TESTL, VERR, VERW, WAIT, WBINVD,
        XADD, XADDB, XADDW, XADDL, XCHG, XCHGB, XCHGW, XCHGL, XLAT, XLATB, XOR, XORB, XORW, XORL},
  keywordstyle=\color{blue}\bfseries,
  ndkeywordstyle=\color{darkgray}\bfseries,
  identifierstyle=\color{black},
  sensitive=false,
  comment=[l]{\#},
  morecomment=[s]{/*}{*/},
  commentstyle=\color{purple}\ttfamily,
  stringstyle=\color{red}\ttfamily,
  morestring=[b]',
  morestring=[b]"
}

\lstdefinelanguage{riscv}{
  keywords={
    LUI, AUIPC,
    JAL, JALR,
    BEQ, BNE, BLT, BGE, BLTU, BGEU,
    LB, LH, LW, LBU, LHU, LD,
    SB, SH, SW, SD,
    ADDI, SLTI, SLTIU, XORI, ORI, ANDI,
    SLLI, SRLI, SRAI,
    ADD, SUB, SLL, SLT, SLTU, XOR, SRL, SRA, OR, AND,
    ADDIW, SLLIW, SRLIW, SRAIW,
    ADDW, SUBW, SLLW, SRLW, SRAW,
    MUL, MULH, MULHSU, MULHU, DIV, DIVU, REM, REMU,
    MULW, DIVW, DIVUW, REMW, REMUW,
    LR.W, SC.W, AMOSWAP.W, AMOADD.W, AMOXOR.W,
    AMOAND.W, AMOOR.W, AMOMIN.W, AMOMAX.W,
    AMOMINU.W, AMOMAXU.W,
    LR.D, SC.D, AMOSWAP.D, AMOADD.D, AMOXOR.D,
    AMOAND.D, AMOOR.D, AMOMIN.D, AMOMAX.D,
    AMOMINU.D, AMOMAXU.D,
    FLW, FSW,
    FMADD.S, FMSUB.S, FNMSUB.S, FNMADD.S,
    FADD.S, FSUB.S, FMUL.S, FDIV.S, FSQRT.S,
    FSGNJ.S, FSGNJN.S, FSGNJX.S,
    FMIN.S, FMAX.S,
    FCVT.W.S, FCVT.WU.S, FCVT.S.W, FCVT.S.WU,
    FMV.X.W, FMV.W.X,
    FEQ.S, FLT.S, FLE.S,
    FLD, FSD,
    FMADD.D, FMSUB.D, FNMSUB.D, FNMADD.D,
    FADD.D, FSUB.D, FMUL.D, FDIV.D, FSQRT.D,
    FSGNJ.D, FSGNJN.D, FSGNJX.D,
    FMIN.D, FMAX.D,
    FCVT.W.D, FCVT.WU.D, FCVT.D.W, FCVT.D.WU,
    FCVT.S.D, FCVT.D.S,
    FMV.X.D, FMV.D.X,
    FEQ.D, FLT.D, FLE.D,
    CSRRW, CSRRS, CSRRC,
    CSRRWI, CSRRSI, CSRRCI,
    FENCE.I,
    FENCE,
    ECALL, EBREAK,
    MRET, SRET, URET,
    WFI,
    SFENCE.VMA,
    NOP, LI, LA, MV,
    NOT, NEG, NEGW,
    SEQZ, SNEZ, SLTZ, SGTZ,
    BEQZ, BNEZ, BLEZ, BGEZ, BLTZ, BGTZ,
    BGT, BLE, BGTU, BLEU,
    J, JR, RET,
    CALL, TAIL,
    RDINSTRET, RDCYCLE, RDTIME,
    CSRR, CSRW, CSRS, CSRC,
    CSRWI, CSRSI, CSRCI
  },
  keywordstyle=\color{blue}\bfseries,
  ndkeywordstyle=\color{darkgray}\bfseries,
  identifierstyle=\color{black},
  sensitive=false,
  comment=[l]{\#},
  morecomment=[s]{/*}{*/},
  commentstyle=\color{purple}\ttfamily,
  stringstyle=\color{red}\ttfamily,
  morestring=[b]',
  morestring=[b]"
}

\lstdefinelanguage{arm}{
  keywords={
    ADC, ADD, AND, BIC, CMN, CMP, EOR, MOV, MVN,
    RSB, RSC, SBC, SUB, TST, TEQ,
    MLA, MLS, MUL, SMLAL, SMULL, UMLAL, UMULL,
    QADD, QSUB, QDADD, QDSUB,
    QADD16, QADD8, QASX, QSAX,
    QSUB16, QSUB8, QSAX, QSUBADDX,
    ASR, LSL, LSR, ROR, RRX,
    B, BL, BX, BLX,
    BXJ,
    BEQ, BNE, BCS, BHS, BCC, BLO,
    BMI, BPL, BVS, BVC,
    BHI, BLS, BGE, BLT, BGT, BLE,
    BAL,
    LDR, LDRB, LDRH, LDRSB, LDRSH,
    LDRD, LDM, LDMIA, LDMIB, LDMDA, LDMDB,
    LDRT, LDRBT, LDRHT, LDRSBT, LDRSHT,
    STR, STRB, STRH, STRD,
    STM, STMIA, STMIB, STMDA, STMDB,
    STRT, STRBT, STRHT,
    PUSH, POP,
    SWP, SWPB,
    REV, REV16, REVSH,
    BFC, BFI, BFXIL, SBFX, UBFX,
    CLZ,
    LDREX, LDREXB, LDREXH, LDREXD,
    STREX, STREXB, STREXH, STREXD,
    DMB, DSB, ISB,
    MCR, MCRR, MRC, MRRC,
    MRS, MSR,
    CPS, SETEND, SVC, BKPT,
    CLREX, WFE, WFI, SEV,
    YIELD, NOP,
    VLDR, VSTR,
    VMOV, VABS, VNEG,
    VADD, VSUB, VMUL, VDIV,
    VMLA, VMLS, VNMLA, VNMLS, VNMUL,
    VMAX, VMIN,
    VSQRT,
    VCMP, VCMPE,
    VCVT,
    VCLZ, VCLS,
    VAND, VORR, VEOR, VBIC, VMVN,
    VLD1, VLD2, VLD3, VLD4,
    VST1, VST2, VST3, VST4,
    ADR, ADCS, ADDS, SUBS, MOVS, ANDS, ORRS, EORS
  },
  keywordstyle=\color{blue}\bfseries,
  ndkeywordstyle=\color{darkgray}\bfseries,
  identifierstyle=\color{black},
  sensitive=false,
  comment=[l]{@},
  morecomment=[l]{;},
  morecomment=[s]{/*}{*/},
  commentstyle=\color{purple}\ttfamily,
  stringstyle=\color{red}\ttfamily,
  morestring=[b]',
  morestring=[b]"
}

\lstset{style=codestyle, language=C}

% disponi sezioni
\usepackage{titlesec}

\titleformat{\section}
	{\sffamily\Large\bfseries} 
	{\thesection}{1em}{} 
\titleformat{\subsection}
	{\sffamily\large\bfseries}   
	{\thesubsection}{1em}{} 
\titleformat{\subsubsection}
	{\sffamily\normalsize\bfseries} 
	{\thesubsubsection}{1em}{}

% tikz
\usepackage{tikz}

% float
\usepackage{float}

% grafici
\usepackage{pgfplots}
\pgfplotsset{width=10cm,compat=1.9}

% disponi alberi
\usepackage{forest}

\forestset{
	rectstyle/.style={
		for tree={rectangle,draw,font=\large\sffamily}
	},
	roundstyle/.style={
		for tree={circle,draw,font=\large}
	}
}

% disponi algoritmi
\usepackage{algorithm}
\usepackage{algorithmic}
\makeatletter
\renewcommand{\ALG@name}{Algorithm}
\makeatother

% disponi numeri di pagina
\usepackage{fancyhdr}
\fancyhf{} 
\fancyfoot[L]{\sffamily{\thepage}}

\makeatletter
\fancyhead[L]{\raisebox{1ex}[0pt][0pt]{\sffamily{\@title \ \@date}}} 
\fancyhead[R]{\raisebox{1ex}[0pt][0pt]{\sffamily{\@author}}}
\makeatother

\begin{document}
% sezione (data)
\section{Lesson 01-10-26}

% stili pagina
\thispagestyle{empty}
\pagestyle{fancy}

% testo
\subsection{Freestanding C}

\lstset{language=c}

The C standard defines two kinds of environment:
\begin{itemize}
	\item \textbf{Hosted}: C library (\texttt{libc}) and an OS sit under the user code, library functions such as \lstinline|prinf()|, \lstinline|malloc()|, \lstinline|fopen()| are available. \lstinline|main()| is called by the \texttt{libc} startup code, called by the OS. Of course, each OS implementation implements its own \texttt{libc} (GNU/Linux, Windows, ...). There are even versions of \texttt{libc} which run on bare metal (\texttt{newlib});
	\item \textbf{Freestanding}: nothing stis under the user code, no library functions are available and nobody calls \lstinline|main()| (the default entry point is the reset handler).
\end{itemize}

We will be using a \textit{freestanding} environment from now on.

\par\smallskip
\textbf{\textsf{System calls}}

In fact, \texttt{libc} in Linux is only an interface between user code and the OS: a call to a \texttt{libc} function in Linux means:
\begin{itemize}
	\item Performing a little setup work in userland, via the \texttt{libc} implementations;
	\item Performing some kind of \textit{system call}, which leaves the CPU to the kernel;
	\item Waiting for the kernel to carry out our system call;
	\item Getting the result back from \texttt{libc}.
\end{itemize}

We decide to opt out of libraries like \texttt{libc} because they still need architecture-specific hooks (\texttt{write}, \texttt{crt0}, ...) hooks to be implemented: we're better off doing everything ourselves.

\par\smallskip
\textbf{\textsf{Compiler helpers}}

One piece of help we will use is the \textit{compiler helpers} provided by the \textbf{GCC} (\textit{Gnu Compiler Collection}) toolchain we use.
These implement some basic C operators (such as divide and modulo), which may not be implemented in hardware.
Was someone to be curious, some example helpers can be found at \url{https://github.com/seggiani-luca/micro-sim/blob/main/eprom/src/lib/helpers.cpp} (there are likely cleaner and better implemented one somewhere in the GCC source).

\subsubsection{Semihosting}

We still need to be able to print stuff to console.
The solution we will use is \textbf{semihosting}: the target asks the \textit{host} (our machine) to perform the I/O.

\begin{enumerate}
\item The program puts an operation code in \texttt{r0} (\texttt{0x04} = \texttt{SYS\_WRITE0}) and the string address in \texttt{r1}. \texttt{SYS\_WRITE0} is the opcode for "print this NULL-terminated string";

\item It executes \texttt{bkpt 0xab}: a breakpoint instruction with a special number;

\item QEMU (or a debugger) sees the breakpoint, reads \texttt{r0} and \texttt{r1}, reads the string from target memory, and prints it \textbf{on your laptop};

\item The program resumes after the \texttt{bkpt}.

\end{enumerate}

This only works with a debugger or emulator attached (e.g., \texttt{-semihosting} in \texttt{make run}). On a real board without a debugger, \texttt{bkpt} would simply stop the CPU.

\par\smallskip
\textbf{\textsf{In code}}

In code, this translates to:
\begin{lstlisting}[language=C++, style=codestyle]	
static void semihost_write0( const char *message )
{
    register uint32_t    r0 __asm__( "r0" ) = 0x04;    /* SYS_WRITE0 */
    register const char *r1 __asm__( "r1" ) = message; /* string address */
    __asm__ volatile ( "bkpt 0xab" : : "r"( r0 ), "r"( r1 ) : "memory" );
}
\end{lstlisting}

Let's look at what \texttt{bkpt 0xab} really does:

\begin{itemize}
\item The program \textbf{stops} doing its own work;
\item Control passes to \textbf{someone more powerful} (QEMU, outside the CPU);
\item That someone performs the job and \textbf{hands control back}.
\end{itemize}

This is the same shape as a system call on Linux: your program cannot drive the terminal itself, so it \textbf{asks the kernel} to do it.

Next question: can the Cortex-M3 itself have a `powerful'' part and a `restricted'' part, like kernel mode and user mode? \textbf{Yes}: privileged and unprivileged execution, next slide.

\subsection{Privilege level}

Even if our implementation of an OS will be \textit{flat}, let's go over \textbf{privilege} levels in our ARM Cortex-M3 CPU.

W know that an \textit{OS kernel} runs in kernel mode, where as an applications runs in \textit{user mode}. Cortex-M3 has the same idea in hardware:
\begin{itemize}
	\item \textbf{Handler mode}: used when running an exception handler, always privileged;
	\item \textbf{Thread mode}: normal mode used for user code, can be privileged or unprivileged.
\end{itemize}

After reset we find the CPU in thread mode, privileged.
We have to manually set the CPU to unprivileged thread mode if we want protection on user code.

A special CPU register (not in memory) selects Thread-mode privilege:

\begin{table}[h]
\centering
\begin{tabular}{c l l}
\hline
\textbf{Bit} & \textbf{Name} & \textbf{Meaning} \\
\hline
0 & \textbf{nPRIV} & 0 = privileged, 1 = unprivileged (Thread mode only) \\
1 & SPSEL & selects which stack pointer Thread mode uses (not used here) \\
\hline
\end{tabular}
\end{table}

Special registers are reached with two dedicated instructions:

\begin{table}[h]
\centering
\begin{tabular}{c l}
\hline
\textbf{Instruction} & \textbf{C equivalent} \\
\hline
\texttt{mrs r0, control} & \texttt{r0 = CONTROL;} \quad (read special register) \\
\texttt{msr control, r0} & \texttt{CONTROL = r0;} \quad (write special register) \\
\hline
\end{tabular}
\end{table}

\subsubsection{Dropping privilege}

So, for example, to drop our privilege lever to unprivileged:
\begin{lstlisting}[language=arm, style=codestyle]	
mrs r0, control     @ r0 = CONTROL
orr r0, r0, #1      @ r0 = r0 | 1      (set nPRIV)
msr control, r0     @ CONTROL = r0     -> now unprivileged
isb                 @ discard instructions already fetched under old rules
\end{lstlisting}

Writing CONTROL is itself a privileged operation: an unprivileged \texttt{msr control}, ... is ignored. Code cannot promote itself back.
The only ways back to privileged: reset, or an exception (its handler runs in Handler mode, always privileged).

\begin{table}[H]
\centering
\begin{tabular}{p{10cm} p{5cm}}
\hline
\textbf{Attempt} & \textbf{Result} \\
\hline
Read/write the \textbf{System Control Space} (\texttt{0xE000E000} \ldots):
SysTick, NVIC, SCB (VTOR, SHCSR, CCR, priorities)
& \textbf{BusFault} \\

\texttt{msr control, r0} (raise privilege back)
& ignored \\

\texttt{cpsid i} / \texttt{cpsie i} (mask interrupts)
& ignored \\

Semihosting \texttt{bkpt 0xab}
& QEMU refuses it by default $\rightarrow$ fault \\
\hline
\end{tabular}
\end{table}

Ordinary RAM, Flash, and user code still work normally. The restricted code simply cannot \textbf{reconfigure} the machine or silence it.

\subsubsection{Raising privilege}

Unprivileged code cannot touch the hardware, but it can still \textit{call} privileged mode handlers: 
\texttt{svc \#n} (\textit{SuperVisor Call}) is allowed in unprivileged code.

\begin{enumerate}
	\item 
It raises the \textbf{SVCall exception} (vector-table entry 11, \texttt{SVC\_Handler}).
The CPU enters \textbf{Handler mode}, and the corresponding handler runs privileged;
	\item
When the handler returns, the CPU goes back to the instruction after \texttt{svc}, still \textbf{unprivileged}.

\end{enumerate}

\texttt{\#n} is a small number encoded in the instruction. Real kernels use a request number and arguments to tell which service is wanted; here we only use \texttt{svc \#0}.

This allows us to implement what effectively are \textbf{system calls}.

\end{document}
